\documentclass[11pt,a4paper]{article}

\usepackage[utf8]{inputenc}
\usepackage[T1]{fontenc}
\usepackage[swedish,english]{babel}
\usepackage[margin=2.5cm]{geometry}
\usepackage{microtype}
\usepackage{xcolor}
\usepackage{tcolorbox}
\tcbuselibrary{breakable}
\usepackage{enumitem}
\usepackage{amsmath,amssymb}
\usepackage{booktabs}
\usepackage{tabularx}
\usepackage{hyperref}
\usepackage{verbatim}

\definecolor{caseblue}{RGB}{0,75,135}
\definecolor{datagreen}{RGB}{39,120,80}

\hypersetup{
    colorlinks=true,
    linkcolor=caseblue,
    urlcolor=caseblue
}

\setlist[itemize]{leftmargin=*,itemsep=0.3em}
\setlist[enumerate]{leftmargin=*,itemsep=0.3em}
\setlength{\emergencystretch}{2em}

\newtcolorbox{casebox}[1]{
    breakable,
    colback=caseblue!5,
    colframe=caseblue,
    title={#1},
    fonttitle=\bfseries
}

\newtcolorbox{databox}[1]{
    breakable,
    colback=datagreen!6,
    colframe=datagreen,
    title={#1},
    fonttitle=\bfseries
}

\title{
    {\Large C1DB1B: Datahantering f\"or beslutsfattande}\\[0.5em]
    {\LARGE Document B: The Data Stream}\\[0.3em]
    {\large \textit{Datastr\"ommen}}\\[0.5em]
    {\large Versioned datasets, dictionaries, and module releases}
}

\author{Course Team}
\date{Current offering \quad $|$ \quad Vintage as stated below}

\begin{document}

\maketitle

\begin{casebox}{Why this document exists}
The assignment (\texttt{project\_Introduction.tex}) and the module handbook
(\texttt{module\_design.tex}) describe \emph{what students decide}
and \emph{how they work}. They should not contain a particular year's
population counts, clinic lists, or SCB table vintages.

This document is the data contract for \textbf{course-issued} files.
When Statistics Sweden, VGR, or the course team updates those files,
only the folders under \texttt{data/00}--\texttt{data/04} and
\texttt{manifest.json} change.

Notebooks \emph{read} the manifest. Students do not edit it.

Group-collected downloads (VGR, SCB, or other public extracts
for the group's approved KPIs) are the expansion of that
architecture. They follow a different path, described below.
They are not optional extras once Module~1 has Verified them.
\end{casebox}

\section{How to Change Data Over Time}

The course uses one folder tree:

\begin{verbatim}
data/
  manifest.json
  00_decision_model_zero/
  01_module1/
  02_module2/
  03_module3/
  04_module4/
\end{verbatim}

\texttt{data/manifest.json} is the single switchboard. It records:
\begin{itemize}
    \item the course vintage (for example \texttt{2026.1});
    \item which file each module must load;
    \item the reference years $t_0$ and $t_1$ used in calculations;
    \item the official source table identifiers and retrieval dates.
\end{itemize}

To refresh a module's data without rewriting teaching materials:
\begin{enumerate}
    \item Place the new extract in the relevant module folder. Keep the
    previous file if you need an audit trail (rename with a date suffix).
    \item Update the corresponding paths, years, and source notes in
    \texttt{manifest.json}.
    \item Update the vintage line at the top of this document and the
    module source note.
    \item Re-run the instructor solution notebook for that module and
    store new reference values if the live challenge depends on them.
\end{enumerate}

Notebooks must load years and paths from the manifest:
\begin{verbatim}
MANIFEST = json.loads(Path("data/manifest.json").read_text(encoding="utf-8"))
t0 = MANIFEST["module1"]["t0"]
t1 = MANIFEST["module1"]["t1"]
anchor = Path(MANIFEST["module1"]["anchor_file"])
bench = Path(MANIFEST["module3"]["benchmark_file"])
\end{verbatim}

Do not write \texttt{2015} or \texttt{2025} as literals in student tasks
unless you are illustrating a historical example.

The numbered steps above are for the \textbf{course team} when a
new vintage is issued. They are not student tasks.

\section{Course Files versus Group-Collected Files}

Students work with two layers of data.

\begin{center}
\small
\begin{tabularx}{\textwidth}{@{}p{0.28\textwidth}p{0.32\textwidth}X@{}}
\toprule
Layer & Who writes the files & How notebooks find them \\
\midrule
Course-issued
(anchor, companion,
benchmark, scenarios) &
The course team. &
Read paths and years from
\texttt{data/manifest.json}.
Do not edit that file. \\[0.8em]
Group-collected
(extra VGR, SCB, or
other public extracts
for approved KPIs) &
The group. &
Save a CSV under
\texttt{data/group\_sources/}
and load it with an explicit
path in the notebook.
Record it in the group's
source register, not in
\texttt{manifest.json}. \\
\bottomrule
\end{tabularx}
\end{center}

\subsection{What students do with a VGR download}

Yes: they format it so it can join the five-municipality table.
No: they do not wire it through \texttt{manifest.json}.

\begin{enumerate}
    \item Keep the original download (for example a PDF table or
    Excel export) outside the issued folders, or in
    \texttt{data/group\_sources/raw/}.
    \item Do not overwrite the downloaded file. In the notebook,
    standardize keys and rename columns to the field names below
    (\texttt{Municipality\_Code}, \texttt{Year}, \ldots).
    \item Save one analysis-ready CSV per source, for example
    \texttt{data/group\_sources/travel\_times.csv}.
    \item Add one row to the group's source register
    (\texttt{data/group\_sources/Source\_Register.csv}):
    file path, observation unit, keys, coverage, units,
    missing-value markers.
    \item Join with the same Module~2 rules:
    \texttt{how}, \texttt{validate}, \texttt{indicator}.
    Keep all five municipalities. A missing match is a coverage
    gap, not a zero.
\end{enumerate}

A CSV is enough. Students are not asked to write JSON.
A broken comma in \texttt{manifest.json} would stop the whole
class pipeline; a broken group CSV affects only that group
and is visible in the source register.

If a group cannot produce a comparable extract for all five
municipalities, the indicator stays Unresolved or Unknown.
That is an acceptable Module~1--2 outcome.

There is no advanced, source-specific support for group
downloads. The course team does not know in advance which
table a group will find. Formatting, name-to-code mapping,
and the join are the group's work, using the issued companions
as the method example.

\begin{databox}{Do not put group files here}
\texttt{data/00\_decision\_model\_zero/} through
\texttt{data/04\_module4/} and \texttt{data/manifest.json}
are course-issued. Replacing them mixes vintages and breaks
the live challenge.
\end{databox}

\begin{databox}{Current vintage}
\textbf{Vintage:} 2026.1\\
\textbf{Released:} with this vintage (not a timetable date)\\
\textbf{Anchor reference years:} $t_0 = 2015$, $t_1 = 2025$\\
\textbf{Geographic coverage:} 52 municipalities present in the SCB
V\"astra G\"otaland extract (including three Halland-border codes used
in the source table: 1382 Falkenberg, 1383 Varberg, 1384 Kungsbacka).\\
\textbf{Predictive component:} not issued as a common file.
Optional for a group: own model (Module~2 may hold out
about 20\% of eligible observations) or an external forecast,
either way justified. Otherwise mark Not assigned.\\
\textbf{Module 3 common exercise:} synthetic WSM/TOPSIS benchmark,
issued.\\
\textbf{Module 4 common exercise:} scenario specification and teaching
cost ceiling, issued. Not a predictive hold-out.
\end{databox}

\section{Document Map}

\begin{center}
\small
\begin{tabularx}{\textwidth}{@{}p{0.22\textwidth}p{0.28\textwidth}X@{}}
\toprule
Document & Role & Changes when? \\
\midrule
\texttt{project\_Introduction.tex} &
The V{\aa}rdlyftet assignment &
Only if the decision question changes. \\[0.6em]
\texttt{module\_design.tex} &
Module tasks, gateways, methods &
Only if pedagogy or assessment changes. \\[0.6em]
\texttt{lectures2026/} &
Preparation lectures &
When teaching sequence or examples change. \\[0.6em]
\texttt{Data\_Stream.tex} (this file)
& Data contract, dictionaries, vintages &
When sources, years, or files change. \\[0.6em]
\texttt{data/manifest.json} &
Machine-readable switchboard &
Every data release. \\
\bottomrule
\end{tabularx}
\end{center}

\section{Release 00: Decision Model Zero}

\textbf{Guiding question from the assignment:}
Which three portfolios would the group prefer to work with, before
systematic methods are taught?

\subsection*{Files}

\begin{itemize}
    \item \texttt{data/00\_decision\_model\_zero/VGRBiddingDataset.csv}
    \item Optional Excel copy:
    \texttt{data/00\_decision\_model\_zero/VGRBiddingDataset.xlsx}
\end{itemize}

\subsection*{What students see}

Municipalities appear as coded keys (\texttt{MUN-VGR-101}, \ldots).
Each row is one municipality inside a five-municipality portfolio
(\texttt{P1}--\texttt{P7}).

\subsection*{Fields}

\begin{center}
\small
\begin{tabularx}{\textwidth}{@{}p{0.28\textwidth}p{0.18\textwidth}X@{}}
\toprule
Field & Type & Meaning \\
\midrule
\texttt{Portfolio ID} & code & Group of five coded municipalities. \\
\texttt{Anonymized Key} & code & Municipality disguise before allocation. \\
\texttt{VGR Region Context} & text & Instructor-only identity (strip before student release if required). \\
\texttt{Total Population (2024/2025)} & integer & Snapshot population. \\
\texttt{Elderly Share \% (65+)} & percent & Broad elderly share, not the Module~1 $80+$ indicator. \\
\texttt{Operating Clinics} & integer & Approximate incumbent count. \\
\texttt{Avg Household Income (SEK/mo)} & SEK & Snapshot income. \\
\texttt{Est. CNI Score} & index & Care Need Index relative to 1.0 national average. \\
\texttt{Primary Trade-off Tension} & text & Qualitative tension, not a scored KPI. \\
\bottomrule
\end{tabularx}
\end{center}

\subsection*{Release rule}

If identities must remain hidden until after allocation, publish a
student copy with \texttt{VGR Region Context} removed. Keep the full
file in the instructor folder.

After allocation, students work with official SCB municipality codes
from Release 01. Do not treat Decision Model Zero figures as the
Module~1 calculation source.

\section{Release 01: Module 1 Anchor}

\textbf{Guiding question from the handbook:}
What should matter in this decision, and can the group calculate a
common demographic baseline for its five municipalities?

The course team provides one regional working table. All groups use it.
Each group then selects the five allocated municipality codes.

\subsection*{Files}

\begin{itemize}
    \item \texttt{data/01\_module1/VGR\_Demographics\_Anchor.csv}
    \item \texttt{data/01\_module1/SOURCE\_NOTE.md}
    \item Path and years in \texttt{data/manifest.json} under
    \texttt{module1}
\end{itemize}

\subsection*{Grain}

One row = one municipality $\times$ one calendar year.

\subsection*{Required fields}

\begin{center}
\small
\begin{tabularx}{\textwidth}{@{}p{0.28\textwidth}p{0.18\textwidth}X@{}}
\toprule
Field & Type & Definition \\
\midrule
\texttt{Municipality\_Code} & identifier & Official SCB kommun-kod. Not a measurement. \\
\texttt{Municipality\_Name} & text & Official name in the extract. \\
\texttt{Year} & integer & Calendar year of the stock count. \\
\texttt{Pop\_Total} & count & Total population, all ages. \\
\texttt{Pop\_0\_17} & count & Ages 0--17. \\
\texttt{Pop\_18\_64} & count & Ages 18--64. \\
\texttt{Pop\_65\_79} & count & Ages 65--79. \\
\texttt{Pop\_80\_Plus} & count & Ages 80 and above, including $100+$. \\
\bottomrule
\end{tabularx}
\end{center}

\subsection*{Provenance (vintage 2026.1)}

\begin{itemize}
    \item \textbf{Producer:} Statistics Sweden (SCB), table BE0101,
    population by region, age, sex, and year.
    \item \textbf{Underlying extract:}
    \texttt{VGR\_Demographics\_2015\_2025.csv}, itself curated from
    \texttt{VGR\_Demographics\_Raw.xlsx}.
    \item \textbf{Construction:} Gender \texttt{totalt} only; age
    \texttt{TOT} excluded; single-year ages aggregated into the four
    bands above; county total (\texttt{region\_code = 14}) excluded.
    \item \textbf{Integrity rule:}
    $\texttt{Pop\_0\_17}+\texttt{Pop\_18\_64}+\texttt{Pop\_65\_79}+\texttt{Pop\_80\_Plus}=\texttt{Pop\_Total}$.
\end{itemize}

\subsection*{Prescribed indicators (formulas stay in the handbook)}

Notebooks must take $t_0$ and $t_1$ from the manifest.

\begin{align*}
S_{80+,t_{1}}
&=
100\,\frac{P_{80+,t_{1}}}{P_{\mathrm{total},t_{1}}},
\\[0.6em]
D_{t_{1}}
&=
100\,\frac{P_{0\text{--}17,t_{1}}+P_{65\text{--}79,t_{1}}+P_{80+,t_{1}}}{P_{18\text{--}64,t_{1}}},
\\[0.6em]
G_{80+}
&=
100\,\frac{P_{80+,t_{1}}-P_{80+,t_{0}}}{P_{80+,t_{0}}}.
\end{align*}

Do not compute mean or median age from these four bands.

\subsection*{Structural checks the notebook must run}

\begin{enumerate}
    \item All five allocated municipality codes are present.
    \item Each municipality has one row for $t_0$ and one row for $t_1$.
    \item \texttt{Municipality\_Code} and \texttt{Year} jointly identify
    rows uniquely.
    \item Required population fields are numeric, non-negative, and not
    missing.
    \item The four age-group counts sum to \texttt{Pop\_Total}.
\end{enumerate}

If a check fails, students inspect extraction or selection. They must
not drop a municipality, replace missing values with zero, or edit the
source file.

\subsection*{Refreshing Module 1}

When SCB publishes a later year, rebuild
\texttt{VGR\_Demographics\_Anchor.csv}, move $t_1$ forward in the
manifest, and keep $t_0$ either fixed (ten-year window) or shifted
(rolling window). Record the choice in the source note.

\section{Release 02: Module 2 Companion Streams}

\textbf{Guiding question from the handbook:}
Can the group combine the sources needed to compare its five
municipalities without losing alternatives, inventing values, or
silently changing indicator meaning?

Module 2 consumes:
\begin{itemize}
    \item the frozen Module~1 anchor for the five allocated
    municipalities (Release 01);
    \item the common companion files in this release;
    \item any additional verified sources from the group's approved
    KPI-to-data table.
\end{itemize}

Complete the common companion exercise first. Then apply the same
checks to portfolio sources. There is no fixed dataset count.

\subsection*{Files}

\begin{itemize}
    \item \texttt{data/02\_module2/VGR\_Companion\_Annual.csv}
    \item \texttt{data/02\_module2/VGR\_Companion\_StaticLookup.csv}
    \item \texttt{data/02\_module2/VGR\_Masked\_TreatmentExercise.csv}
    \item \texttt{data/02\_module2/SOURCE\_NOTE.md}
    \item Paths and status in \texttt{data/manifest.json} under
    \texttt{module2}
\end{itemize}

\subsection*{Predictive component (vintage 2026.1)}

Not issued as a common file. Prediction is optional for a group:
build a model (Module~2 may hold out about 20\% of eligible
observations) or join a published forecast, and justify either.
Do not split the five-municipality decision table as an MCDA
training/test set. Do not split the ten-row Module~1 prototype
to create a hold-out file. If the group does not take it up,
mark \texttt{Not assigned}.

\subsection*{A. Annual companion}

\textbf{Grain:} one row = one municipality in year $t_1$
(the year is stored in the file and must match
\texttt{module1.t1} in the manifest).

\textbf{Join:} \texttt{Municipality\_Code} (and \texttt{Year}, after
confirming both tables use the same $t_1$). Expected relationship
to the $t_1$ demographic extract: one-to-one.

\begin{center}
\small
\begin{tabularx}{\textwidth}{@{}p{0.30\textwidth}p{0.16\textwidth}X@{}}
\toprule
Field & Type & Definition \\
\midrule
\texttt{Municipality\_Code} & identifier & Official SCB kommun-kod. \\
\texttt{Municipality\_Name} & text & Display label only. \\
\texttt{Year} & integer & Calendar year of the observation ($t_1$). \\
\texttt{Operating\_Clinics} & count & Course-constructed incumbent
count. Not a live VGR clinic register. \\
\texttt{Clinics\_per\_10k} & rate &
$10000 \times \texttt{Operating\_Clinics} / \texttt{Pop\_Total}$
at $t_1$. Verify; do not treat as an independent measurement. \\
\texttt{Avg\_Travel\_Time\_min} & minutes & Curated accessibility
estimate. Documented missing marker: \texttt{..}. \\
\bottomrule
\end{tabularx}
\end{center}

Documented travel-time gaps in vintage 2026.1:
\texttt{1444} Gr\"astorp, \texttt{1446} Karlsborg, \texttt{1447}
Gullsp\aa ng, \texttt{1465} Svenljunga.

\subsection*{B. Static municipality lookup}

\textbf{Grain:} one row = one municipality. There is \textbf{no}
\texttt{Year} field.

\textbf{Join:} \texttt{Municipality\_Code} only. After the annual
table has been reduced to one row per municipality, the expected
relationship is one-to-one. If annual history is joined before
filtering to $t_1$, the relationship is many-to-one.

Do not invent a year column on this file. If a static value is reused
at $t_1$, label the reference period and the assumption.

\begin{center}
\small
\begin{tabularx}{\textwidth}{@{}p{0.30\textwidth}p{0.16\textwidth}X@{}}
\toprule
Field & Type & Definition \\
\midrule
\texttt{Municipality\_Code} & identifier & Official SCB kommun-kod. \\
\texttt{Municipality\_Name} & text & Display label only. \\
\texttt{Care\_Need\_Index} & index & Course-constructed CNI relative
to 1.0. Teaching values aligned with Decision Model Zero where
identities match; otherwise 1.0. \\
\texttt{Est\_Premises\_Cost\_kSEK} & kSEK & Course-constructed
premises-cost teaching estimate. Documented missing marker:
\texttt{..} (vintage 2026.1: \texttt{1438} Dals-Ed). \\
\bottomrule
\end{tabularx}
\end{center}

\subsection*{C. Masked treatment exercise}

\textbf{Grain:} one synthetic record. These rows are \textbf{not}
municipalities and must not enter the decision table.

\begin{center}
\small
\begin{tabularx}{\textwidth}{@{}p{0.28\textwidth}p{0.18\textwidth}X@{}}
\toprule
Field & Type & Definition \\
\midrule
\texttt{Record\_ID} & identifier & Teaching record label. \\
\texttt{Synthetic\_Value} & numeric or \texttt{..} & Observed values
10, 12, 14, 64 and one documented missing marker. \\
\texttt{Note} & text & \texttt{observed}, \texttt{observed outlier},
or \texttt{documented missing}. \\
\bottomrule
\end{tabularx}
\end{center}

Use this file for the required two-treatment comparison when the
portfolio has no suitable missingness issue, or as the common
mean-versus-median rehearsal. Keep results separate from factual
establishment evidence.

\subsection*{Missing-value rule}

The documented marker in this release is \texttt{..}.
Parse it as missing. It is not zero, not ``no clinics'', and not
``zero minutes''.

Do not treat every blank, dash, or unmatched join as the same event:
\begin{itemize}
    \item missing in the source;
    \item unmatched after a join;
    \item not applicable or suppressed, where documented;
    \item unavailable because coverage differs.
\end{itemize}

\subsection*{Required pipeline checks}

\begin{enumerate}
    \item Every input has an observation unit, key fields, coverage,
    units, and missing-value markers in the source register.
    \item Keys are standardized in code. Source files remain unmodified.
    \item Key uniqueness is checked before each join.
    \item Joins declare \texttt{how}, \texttt{validate}, and
    \texttt{indicator}.
    \item The audit records input counts, output count, unmatched keys,
    and the join relationship.
    \item All five allocated codes remain in the decision table.
    \item A missing match is reported, not filled with zero and not
    used as a reason to drop a municipality.
    \item Predictive partitions are omitted unless the group took
    up a model in Module~2 (then hold out about 20\% of eligible
    observations). An external forecast is joined like any other
    group source. Do not split the five municipalities.
\end{enumerate}

\subsection*{What these files are not}

They are not a live clinic register, a cadastral premises-cost extract,
or an official Care Need Index publication. Labelled teaching fields
must not be presented as independently verified establishment facts.
Group-collected external sources remain the group's responsibility
and must pass the same checks.

\subsection*{Refreshing Module 2}

Replace files only inside \texttt{data/02\_module2/}. Update
\texttt{manifest.json} paths, the predictive-component flag, and this
section. Re-execute the instructor solution notebook. If a later
vintage assigns prediction, publish the modelling specification
before requiring a split.

\section{Release 03: Module 3 Benchmark}

\textbf{Guiding question from the handbook:}
Under explicit preferences, which alternative performs best in a
checked comparison, and what would change that conclusion?

Module 3 consumes:
\begin{itemize}
    \item the group's approved five-municipality decision table
    (produced in Module~2 from Releases 01 and 02);
    \item the common synthetic benchmark in this release;
    \item the group's recorded indicator directions and baseline
    weights.
\end{itemize}

The course team does \textbf{not} issue a new municipal extract for
Module~3. The benchmark is a calculation exercise. Portfolio scores
are computed from the Module~2 table.

Complete and check the benchmark in \texttt{starter\_m3.ipynb} before
applying the same functions to the portfolio.

\subsection*{Files}

\begin{itemize}
    \item \texttt{data/03\_module3/MCDA\_Benchmark.csv}
    \item \texttt{data/03\_module3/MCDA\_Benchmark\_Meta.csv}
    \item \texttt{data/03\_module3/SOURCE\_NOTE.md}
    \item Paths and status in \texttt{data/manifest.json} under
    \texttt{module3}
\end{itemize}

\subsection*{Predictive component (vintage 2026.1)}

Not issued as a common file. Insert a predicted column only if
the group justified it in Module~2 and the value is finite for
all five municipalities. Label it as a prediction. Otherwise mark
\texttt{Not assigned}. Do not invent a regression target merely
to fill the matrix.

\subsection*{A. Benchmark matrix}

\textbf{Grain:} one row = one synthetic alternative (\texttt{A}--\texttt{E}).
These labels are not municipalities and must not enter the portfolio
decision table.

\begin{center}
\small
\begin{tabularx}{\textwidth}{@{}p{0.28\textwidth}p{0.16\textwidth}X@{}}
\toprule
Field & Type & Definition \\
\midrule
\texttt{Alternative} & identifier & Teaching label A--E. \\
\texttt{opportunity} & index & Benefit criterion; higher preferred. \\
\texttt{annual\_cost} & SEK million & Cost criterion; lower preferred. \\
\bottomrule
\end{tabularx}
\end{center}

Values (vintage 2026.1):

\begin{center}
\small
\begin{tabular}{lrr}
\toprule
Alternative & opportunity & annual\_cost \\
\midrule
A & 10 & 10 \\
B & 8 & 6 \\
C & 6 & 4 \\
D & 4 & 3 \\
E & 2 & 2 \\
\bottomrule
\end{tabular}
\end{center}

\subsection*{B. Benchmark metadata}

\textbf{Grain:} one row = one scored indicator.

\begin{center}
\small
\begin{tabularx}{\textwidth}{@{}p{0.22\textwidth}p{0.16\textwidth}X@{}}
\toprule
Field & Type & Definition \\
\midrule
\texttt{indicator} & name & Must match a column in the benchmark
matrix. Align weights and directions by this name, not by column
order. \\
\texttt{direction} & \texttt{Benefit} or \texttt{Cost} &
Benefit: higher preferred. Cost: lower preferred. \\
\texttt{weight} & $[0,1]$ & Non-negative. The issued baseline sums
to 1: opportunity $0.60$, annual cost $0.40$. \\
\bottomrule
\end{tabularx}
\end{center}

\subsection*{Method conventions (lecture, notebook, and this contract)}

\begin{itemize}
    \item WSM uses direction-aware min--max normalization. A constant
    column is scored 0 for every alternative; retain its weight.
    \item TOPSIS vector-normalizes the \textbf{raw} matrix, then applies
    weights. Do not feed the WSM $Z$ matrix into TOPSIS.
    \item Cost criteria keep non-negative weights. Direction is handled
    in normalization, not by a negative $w$.
    \item Rank higher scores first. Exact ties use average rank.
    \item Spearman $\rho$ is descriptive. It is undefined if a ranking
    is constant. No significance test is required for five alternatives.
\end{itemize}

\subsection*{Reference checks (vintage 2026.1)}

WSM scores: A $0.60$, B $0.65$, C $0.60$, D $0.50$, E $0.40$.

TOPSIS closeness (three decimal places): A $0.565$, B $0.648$,
C $0.588$, D $0.487$, E $0.435$.

Live-challenge rehearsal on this same file: set annual-cost weight
to $0.70$ and redistribute. WSM then ranks E first with score $0.70$.
The in-class instruction may use a different $\alpha$; it still uses
this benchmark, not an unavailable portfolio indicator.

\subsection*{Portfolio scoring}

A Module~2 column may enter the scoring matrix only if:
\begin{itemize}
    \item it is an approved scored indicator;
    \item every one of the five municipalities has a finite value;
    \item Benefit/Cost direction is recorded;
    \item it is not an identifier, year, or status flag.
\end{itemize}

Exclude a source-missing field (vintage 2026.1 example: travel time
for \texttt{1465} Svenljunga) unless an instructor-approved treatment
exists. Do not fill with zero to force a rank. Do not drop a
municipality to make a method run.

Clinic counts, CNI, and premises costs remain labelled teaching
values. They may be scored only if they were approved in Module~1,
and the label stays in the metadata.

\subsection*{What these files are not}

They are not V\"astra G\"otaland municipalities, not a live cost
extract, and not establishment evidence. Benchmark leaders are
teaching results.

\subsection*{Refreshing Module 3}

Replace files only inside \texttt{data/03\_module3/}. Update
\texttt{manifest.json} paths and the predictive-component flag.
Re-execute \texttt{solution\_m3.ipynb} and replace the reference
scores in this section and in the source note if the benchmark
values change.

\section{Release 04}

\texttt{data/04\_module4/} issues the Module~4 scenario contract.
It is not a predictive hold-out.

\subsection{Files}

\begin{center}
\small
\begin{tabularx}{\textwidth}{@{}p{0.34\textwidth}X@{}}
\toprule
File & Role \\
\midrule
\texttt{Scenario\_Spec.csv} &
Named changes for baseline, operating, and combined scenarios
on the synthetic benchmark and the portfolio cost field. \\[0.4em]
\texttt{Constraint\_Spec.csv} &
Teaching ceilings. The benchmark ceiling is compulsory for the
common rehearsal. The portfolio ceiling is an illustrative
demonstration assumption. \\[0.4em]
\texttt{SOURCE\_NOTE.md} &
Source note and live-challenge boundary. \\
\bottomrule
\end{tabularx}
\end{center}

\subsection{Issued parameters}

\begin{itemize}
    \item Operating scenario: multiply the named cost field by $1.20$.
    Keep other indicators and baseline weights unchanged.
    \item Combined scenario (benchmark): the same cost scaling and
    cost weight $0.70$ (opportunity $0.30$).
    \item Combined scenario (portfolio): the same cost scaling and
    the Module~3 premises-cost priority weight $0.40$.
    \item Benchmark ceiling: \texttt{annual\_cost} $\le 6.5$ cost units.
    Synthetic teaching assumption. Alternative B ($6.0$) meets it at
    baseline and exceeds it after the 20\% increase ($7.2$).
    Alternative A ($10$) already exceeds it at baseline.
    \item Demonstration portfolio ceiling:
    \texttt{Est\_Premises\_Cost\_kSEK} $\le 650$.
    Illustrative teaching assumption for instructor notebooks.
    Groups must label any portfolio ceiling as a group assumption
    unless a case requirement is issued.
\end{itemize}

\subsection{What this release is not}

Do not split the five-municipality table.
Do not treat these files as reserved outcomes.
Prediction remains \textbf{Not assigned} in vintage 2026.1.

Uniform positive scaling of one criterion leaves prescribed
normalized values unchanged when normalization is recomputed.
Absolute costs still rise. Recheck the ceiling.

To refresh this release, replace files only inside
\texttt{data/04\_module4/}, update \texttt{manifest.json},
re-execute \texttt{solution\_m4.ipynb}, and revise the reference
notes in this section if the issued multiplier or ceiling changes.

\section{Field Naming Convention}

Use the following names in every new extract so notebooks remain
stable:
\begin{itemize}
    \item \texttt{Municipality\_Code} --- SCB kommun-kod, stored as text
    or integer identifier, never used as a numeric feature.
    \item \texttt{Municipality\_Name} --- display label only.
    \item \texttt{Year} --- calendar year of the observation.
    \item \texttt{Alternative} --- synthetic benchmark row label in
    Module~3 only.
    \item \texttt{direction} --- \texttt{Benefit} or \texttt{Cost}.
    \item \texttt{weight} --- non-negative preference weight.
    \item Indicator-specific fields documented in the module source
    note, not invented ad hoc in notebooks.
\end{itemize}

If an upstream table uses Swedish headers, map them in the source note.
Do not require students to guess column names from a live website.

\section{Instructor Checklist for a Data Refresh}

\begin{enumerate}
    \item Confirm the assignment question in \texttt{project\_Introduction.tex}
    is unchanged. If the decision changed, that is a course change, not
    a data change.
    \item Replace files only inside \texttt{data/<module>/}.
    \item Update \texttt{manifest.json} vintage, dates, paths, and
    $t_0$/$t_1$.
    \item Update this document's vintage box and the relevant source
    note.
    \item Execute the instructor solution notebook
    (\texttt{solution\_m1.ipynb}, \texttt{solution\_m2.ipynb},
    \texttt{solution\_m3.ipynb}, or \texttt{solution\_m4.ipynb})
    and archive reference outputs.
    \item Tell students which vintage they must download. Mixed vintages
    in one group are not permitted.
\end{enumerate}

\begin{casebox}{The data stream in one sentence}
Keep the V{\aa}rdlyftet decision fixed, keep the methods in the handbook
fixed, and let \texttt{data/} plus this document absorb every change in
years, tables, and extracts.
\end{casebox}

\end{document}
